%=============================================================================
\section{Inference engine: scenario inference under partial observability}
\label{sec:framework}
%=============================================================================

\paragraph{Positioning relative to existing approaches.}
The engine separates two tasks. Variables measured in the national survey are
generated from national data alone; outcomes the survey does not measure are
attached through an explicit transfer model built from published aggregates. The objective is decision-relevant inference
under data constraints, not optimal estimation under a fully specified
generative model. Deep generative models such as CTGAN \citep{xu2019} and
variational autoencoders \citep{ma2020} reach high distributional fidelity but,
as commonly implemented, learn from a single table and offer no explicit route
for aggregate external evidence. Bayesian-network generators
\citep{zhang2017,martins2024} provide an interpretable dependency structure but
have likewise been applied to one dataset at a time. BVS keeps the
Bayesian-network generator of \citet{martins2024} unchanged for observed
variables and adds a transfer component for unobserved ones, so that each part
can be checked on its own terms.
Figure~\ref{fig:pipeline_steps} shows the five steps.

\begin{figure}[H]
\centering
\includegraphics[width=0.93\textwidth]{fig_pipeline_steps.png}
\caption{BVS pipeline. Observed variables are learned and generated from
national data alone (Steps~1, 3 and~4). External evidence enters only the
outcome the national survey does not measure, through the transfer model
(Step~2), which later national data update by conjugate fusion.}
\label{fig:pipeline_steps}
\end{figure}

%-----------------------------------------------------------------------------
\subsection{Step 1: Learning the dependency structure from national data}
\label{sec:step1}
%-----------------------------------------------------------------------------

Structure learning identifies which CIS variables are statistically dependent
on which others, and quantifies structural uncertainty instead of committing to
one graph. As \citet{martins2024} show, selecting a single ``optimal'' network
discards the uncertainty inherent in structure learning and can produce worse
synthetic data than averaging over plausible structures.

We estimate a posterior distribution over directed acyclic graphs with a
constrained Bayesian-network sampler. To scale to 90 variables, variables are
grouped into 15 thematic blocks aligned with the ISSP questionnaire, and the
sampler explores each block independently; cross-block dependencies are
recovered through a separate four-layer screening procedure. Domain knowledge
enters as forbidden edges \citep{heckerman1995} that impose an ordering taken
from environmental psychology \citep{ajzen1991,stern2000}: demographic
variables have no parents among attitudes, and general attitudes precede
specific behaviours. This ordering constrains edge directions. It does not turn
observational associations into causal effects, and the learned graph is read
throughout as a dependency structure (Figure~\ref{fig:dag}). Algorithmic
details are given in Appendix~\ref{app:technical}.

\begin{figure}[htbp]
\centering
\includegraphics[width=0.82\textwidth]{fig_dag.png}
\caption{Dependency structure learned from CIS microdata, summarised by
thematic block (seed 42). Line width is the number of edges between two blocks;
outlined in orange are the blocks containing the five predictors of the
transfer model. The target outcomes are not nodes of the network.}
\label{fig:dag}
\end{figure}

%-----------------------------------------------------------------------------
\subsection{Step 2: External evidence for the target outcomes}
\label{sec:step2}
%-----------------------------------------------------------------------------

For each target construct, the category proportions reported for each matching
source item are converted into Dirichlet pseudo-counts,
\begin{equation}
\alpha_{k,s} = n_s \hat{\pi}_{k,s},
\label{eq:pseudocounts}
\end{equation}
where $n_s$ is the sample size and $\hat{\pi}_{k,s}$ the proportion in category
$k$ for entry~$s$. When source and target scales differ (a mean on a
seven-point scale, or a single percentage, mapped to a five-category target),
proportions are redistributed over the target categories by parametric
discretisation. Entries for the same construct are summed
(Figure~\ref{fig:step2_pooling}).

\paragraph{Pooling and heterogeneity.}
Heterogeneity across entries is quantified with the $I^2$ statistic
\citep{higgins2002} and reported in Table~\ref{tab:i2_summary}; it is not
estimable for constructs with a single entry. In the reported specification
$I^2$ does not discount the prior ($\lambda = 0$): the effective sample size
equals the pooled $n$. Two consequences must be stated plainly. First, summing
pseudo-counts treats all entries as draws from one distribution. For a category
with pooled mean $\mu_k$ and total concentration $A$, the prior variance is
$\mu_k(1-\mu_k)/(A+1)$, so adding entries narrows the prior; disagreement
between studies does not widen it. Between-study heterogeneity is therefore not
represented in the intervals. Second, several entries often come from the same
respondents (different items of one questionnaire), so the pooled $n$ exceeds
the number of independent respondents. Table~\ref{tab:i2_summary} reports both.
For perceived safety, for example, six entries from one sample of 116
respondents give a pooled $n$ of 696. The pooled prior
should be read as a transportable summary of the available evidence whose
stated precision is an upper bound.

\begin{table}[htbp]
\centering
\caption{External evidence for the target outcomes. \emph{Entries}: extracted
item-level distributions. \emph{Pooled $n$}: sum of entry sample sizes, equal
to the prior effective sample size. \emph{Indep.\ $n$}: respondents counted
once per study. $I^2$ is not estimable (n.e.) with a single entry.}
\label{tab:i2_summary}
\small
\begin{tabular}{lrrrrr}
\toprule
\textbf{Construct} & \textbf{Studies} & \textbf{Entries} &
  \textbf{Pooled $n$} & \textbf{Indep.\ $n$} & $I^2$ (\%) \\
\midrule
%%I2ROWS%%
\bottomrule
\end{tabular}
\end{table}

\paragraph{Outcomes absent from the national data.}
For the target outcomes no conditional distribution can be learned from the
CIS. For food-packaging acceptance, the outcome with a published base rate and
published profile effects, the conditional probability is built with a
transfer model,
\begin{equation}
P(Y = 1 \mid \mathbf{x}) = \Phi\!\Bigl(c + a_0 \sum_j \beta_j x_j\Bigr),
\qquad
\frac{\sum_{i} w_i\, \Phi\!\bigl(c + a_0 \sum_j \beta_j x_{ij}\bigr)}{\sum_i w_i} = p_0 ,
\label{eq:transfer}
\end{equation}
where $\Phi$ is the standard normal distribution function, $w_i$ the CIS
final weights, $\beta_j$ are
coefficients of an ordered-probit model estimated in the source population,
$a_0 = 0.80$ attenuates them for transfer, and the intercept $c$ is calibrated
so that the weighted average over the $n$ CIS respondents equals the base rate
$p_0$.
The probit link matches the source model, so the coefficients are used on the
scale on which they were estimated. Table~\ref{tab:coefficients} lists them.
The other target outcomes have no published profile effects of comparable
quality and are not projected (Section~\ref{sec:evidence_base}).

\begin{table}[htbp]
\centering
\caption{Inputs of the transfer model for food-packaging acceptance. All
coefficients are from model~2 of \citet{ruokamo2022} (ordered probit,
Finland, $n = 272$), which contains sociodemographic characteristics and
plastic-related practices and attitudes. Effects not reported as significant
in that model are set to zero.}
\label{tab:coefficients}
\small
\begin{tabular}{llrr}
\toprule
\textbf{Predictor} & \textbf{Coding in CIS~3391} & $\beta_j$ & \textbf{SE} \\
\midrule
Female & SEX $= 2$ & 0.401 & 0.148 \\
Under 30 & $2023 - \text{BIRTH} < 30$ & 0.456 & 0.255 \\
Higher education & ESTUDIOS $= 7$ & 0.606 & 0.332 \\
Recycles always or often & V52 $\in \{1, 2\}$ & 0.398 & 0.181 \\
Environmental worry & V15, centred, scaled to $[-1, 1]$ & 0.772 & 0.154 \\
\midrule
Base rate $p_0$ & suitable for food packaging ($n = 301$) & 0.38 & \\
\bottomrule
\end{tabular}
\end{table}

Four assumptions of Eq.~\eqref{eq:transfer} cannot be checked with the
available data and are varied in sensitivity analyses. First, the level: $p_0$
is the share of a Finnish sample that judged recycled plastic suitable for food
packaging, and it is assumed to hold for the Spanish population as a whole.
The transfer factor does not correct it. Second, the effects: sign and ordering
of the coefficients are assumed stable across populations, and $a_0$ shrinks
their size. No Spanish data identify $a_0$ for the target outcome, but its
order of magnitude can be checked on outcomes measured in both countries
(Section~\ref{sec:transport}): Spanish coefficients are %%TA0%% (95\%
interval %%TA0CI%%) times the Finnish ones, so the central value $0.80$ lies
within the empirical range. All results are also reported over $0.50$--$1.00$,
for the two profiles and for every segment (Appendix~\ref{app:sensitivity}). Third, the outcome: the coefficients
were estimated for agreement that recycled content makes a product more
attractive, and are applied to suitability for food packaging, whose base rate
comes from a different item of the same survey. No source measures both items
on the same respondents, so the mapping cannot be tested with the available
data; it carries the direction and ordering of the effects, while the level is
fixed by $p_0$ alone. The projection is accordingly a scenario for suitability
under the assumption that the drivers of attractiveness and of suitability
coincide in sign and ordering. Fourth, the scale of environmental worry: the
five-point item (``not at all'' to ``extremely'' worried) is centred on its
midpoint and scaled to $[-1, 1]$, so the coefficient is the shift from the
midpoint to either extreme. Coding the coefficient as the full contrast between
the extremes gives a spread of 67~points; applying it per scale point gives
81~points. Uncertainty is propagated by Monte Carlo
(5{,}000 draws) over the published standard errors of the coefficients and the
binomial error of $p_0$; cases with missing predictors are excluded
($n = 2{,}133$ of $2{,}254$).

\paragraph{Empirical transport error.}
\label{sec:transport}
The assumptions on level and effects can be confronted with data on outcomes
other than the target. CIS~3391 is the Spanish component of ISSP~2020
Environment~IV \citep{issp2023}, which asks the same questions in %%TNC%%
countries, Finland among them. For five proxy outcomes asked in every country
(willingness to pay higher prices, to pay higher taxes, and to cut one's
standard of living to protect the environment; doing what is right for the
environment despite cost; avoiding products for environmental reasons), we fit
in each country the probit of Eq.~\eqref{eq:transfer} with the same five
predictors and coding, using the ISSP weights (\texttt{transport\_issp.py}).
Three quantities follow. The slope of Spanish on Finnish coefficients, through
the origin and pooled over outcomes and predictors, is %%TA0%%
(%%TA0CI%%, respondent bootstrap): effects transfer from Finland to Spain with
moderate attenuation. The residual Spanish--Finnish discrepancy per
coefficient, beyond sampling error, has a standard deviation of %%TETA%% on
the probit scale; it is largest for gender, whose effect is on average
%%TFEMALE%% lower in Spain. Levels do not transfer: the Spanish share differs
from the Finnish one by a root mean square of %%TLEVEL%% on the probit scale,
and across all %%TNC%% countries coefficients vary with a median
between-country standard deviation of %%TTAU%% times their mean. These
quantities are measured on proxies, not on the target outcome; they set the
order of magnitude of transport error, and \texttt{transfer\_transport.py}
propagates them into the projections (Section~\ref{sec:sensitivity}).

\begin{figure}[h!]
\centering
\includegraphics[width=0.75\textwidth]{fig_step2_pooling.png}
\caption{External evidence after the review rules. Four of the 11 studies
supply retained entries for a target outcome. Only food-packaging acceptance
is projected, through the transfer model; the pooled distributions of the
other outcomes are reported as evidence (Tables~\ref{tab:i2_summary}
and~\ref{tab:novel}) and do not enter observed variables.}
\label{fig:step2_pooling}
\end{figure}

%-----------------------------------------------------------------------------
\subsection{Step 3: Parameter estimation and the role of external evidence}
\label{sec:step3}
%-----------------------------------------------------------------------------

\paragraph{Observed variables.}
For each variable $X_j$ observed in the CIS and each parent configuration $c$,
the conditional probability table is estimated from weighted national counts
$\tilde{\mathbf{n}}_{j,c}$ under a BDeu prior with equivalent sample size 10,
as in \citet{martins2024}:
\begin{equation}
\boldsymbol{\theta}_{j,c} \sim \mathrm{Dir}\!\left(\frac{10}{K_j q_j}\,\mathbf{1} + \tilde{\mathbf{n}}_{j,c}\right),
\label{eq:cpt}
\end{equation}
where $K_j$ is the number of categories of $X_j$, $q_j$ the number of parent
configurations, and the weights are normalised to sum to the sample size.
No external evidence enters these tables. An earlier configuration added
pooled external pseudo-counts to the tables of eleven observed variables. It
was abandoned for three reasons. First, it departs from the single-dataset
estimation for which the method of \citet{martins2024} was developed and
validated. Second, the external items are not equivalent to the CIS items they
were mapped to: several measure different constructs (weekly meat consumption
was treated as reduced plastic consumption, interpersonal trust as institutional
trust), others use reversed response scales, and demographic priors encoded the
sample composition of foreign convenience samples. Third, the evidence was
negative where it could be tested: recovering masked CIS variables from the
external priors alone gave a mean total variation distance of $0.110$ against
$0.038$ with a flat prior, and the external priors did not improve fidelity at
any subsample size from 100 to 1{,}000 respondents (deposit v1.0,
\texttt{level4\_heldout\_summary.json} and \texttt{subsample\_experiment.csv}).

\paragraph{Target outcome and later updating.}
Food-packaging acceptance is attached to every synthetic individual through
Eq.~\eqref{eq:transfer}, evaluated at that individual's own characteristics.
The transfer model then serves as the prior for the first Spanish measurement.
Writing its projection for a segment as a Beta distribution with mean $\pi$ and
effective sample size $m$, a Spanish survey with $y$ acceptances among $N$
respondents of that segment gives
\begin{equation}
\pi \mid y \sim \mathrm{Beta}\bigl(m\pi + y,\; m(1-\pi) + N - y\bigr),
\label{eq:fusion}
\end{equation}
so national data progressively replace the transferred evidence without
re-estimating the network. The prior is obtained by matching a Beta
distribution to the Monte Carlo draws of the projection for the segment,
$m = \bar{p}(1-\bar{p})/s^2 - 1$, where $\bar{p}$ and $s^2$ are the mean and
variance of the draws. The concentration $m$ is therefore not chosen; it
expresses, in respondents, the precision that the transfer inputs imply. Under
the central specification it is 26 for the most receptive profile of
Section~\ref{sec:spread} and 80 for the least receptive, so a first Spanish
sample of 100 respondents of either profile would outweigh the transferred
evidence. Because the intervals omit transfer and extraction uncertainty
(Section~\ref{sec:step5}), these values overstate the weight the transfer
deserves. With a single source study, a between-study variance component cannot
be estimated, so the discount is applied as a power prior \citep{ibrahim2000}:
the prior enters Eq.~\eqref{eq:fusion} with concentration $\delta m$,
$0 < \delta \leq 1$. The empirical transport error of
Section~\ref{sec:transport} gives a direct value: propagated into the draws,
it lowers $m$ to between %%TMLOW%% and %%TMHIGH%% across segments, between a
tenth and a third of the values from parameter uncertainty alone. The measurement plan
therefore uses the transport-based prior (Table~\ref{tab:plan}). The update is implemented
(\texttt{transfer\_query.py}, \texttt{measurement\_plan.py}) and has not been
applied to Spanish data. Conjugate fusion is thus the updating mechanism of
the system, applied where national data are missing, not a correction of
variables the national survey already measures.

%-----------------------------------------------------------------------------
\subsection{Step 4: Posterior predictive generation}
\label{sec:step4}
%-----------------------------------------------------------------------------

A synthetic population is generated by ancestral sampling from the network of
Step~3. To propagate structural uncertainty, $M = 50$ graphs are drawn from the
posterior of Step~1 and 200 individuals are generated from each, giving
10{,}000 synthetic profiles (Figure~\ref{fig:step4}). Food-packaging
acceptance is then drawn for each individual from Eq.~\eqref{eq:transfer}, with
one Monte Carlo draw of the coefficients and of $p_0$ per graph, so that
parameter uncertainty is carried into the population. The age predictor
(under 30) is not identified by the five birth-year bins of the network; it is
drawn from the share of respondents under 30 within each bin in the CIS.
Simulated profiles reduce Monte Carlo error only; they are not
new observations and add no information beyond the evidence already in the
model.

\begin{figure}[htbp]
\centering
\includegraphics[width=0.95\textwidth]{step4.png}
\caption{Synthetic population. Observed variables are generated from the
network; food-packaging acceptance is attached from Eq.~\eqref{eq:transfer}.}
\label{fig:step4}
\end{figure}

%-----------------------------------------------------------------------------
\subsection{Step 5: Validation design and uncertainty accounting}
\label{sec:step5}
%-----------------------------------------------------------------------------

Validation is organised by what each test can and cannot establish
\citep{gelman2013}.
\begin{enumerate}
\item \emph{Reproduction of observed variables}: distributional fidelity and
  dependency preservation of the synthetic population against CIS microdata,
  with baseline generators for comparison.
\item \emph{External comparison}: constructs shared with an independent Spanish
  survey are compared with the synthetic population.
\item \emph{Target outcome}: food-packaging acceptance cannot be validated
  until it is measured in Spain. The first two levels test the generator; they
  do not confirm the projection.
\end{enumerate}

Three kinds of query are computed differently. \emph{Fully specified
profiles}: the probability follows directly from Eq.~\eqref{eq:transfer}, and
its interval from 5{,}000 Monte Carlo draws of the coefficients and of $p_0$;
the network does not enter. \emph{Segments defined by the five predictors}
(Table~\ref{tab:segments}): weighted averages of Eq.~\eqref{eq:transfer} over
the CIS respondents in the segment, reported at the central specification.
\emph{Queries that leave characteristics unspecified or involve other observed
variables} (for instance, institutional trust): these need the joint
distribution of the unspecified characteristics, which the synthetic
population supplies. This is the role of the network. Each individual carries
one draw of the transfer inputs per graph, so these queries combine parameter
and structural uncertainty, and they inherit the cross-block attenuation of
Section~\ref{sec:res_fidelity}.

The reported intervals reflect the published uncertainty of the transfer
inputs. They do not include
uncertainty from extraction and scale harmonisation, from cross-national
transfer, or from between-study heterogeneity. They are therefore lower bounds
on the uncertainty of the projections.

